\chapter{Wnioski}
\label{cha:wnioski}

\section{Podsumowanie realizacji}
\label{sec:podsumowanie}

W ramach projektu zrealizowano kompletny pipeline klasyfikacji znaków dłoni:
od przygotowania danych tablicowych, przez uczenie i ewaluację modeli,
po inferencję z kamery w czasie rzeczywistym. Implementacja zachowuje
spójny format wejścia (28$\times$28, 784 cechy), co pozwala stosować
różne klasyfikatory bez zmian interfejsu danych.

\section{Ocena zastosowanych metod}
\label{sec:ocena_metod}

Przeprowadzone porównanie pokazuje, że:
\begin{itemize}
    \item modele klasyczne (\texttt{SVC}, \texttt{RandomForest}) zapewniają
          zadowalającą bazową skuteczność,
    \item model konwolucyjny (\texttt{CNN}) uzyskuje wyraźnie wyższą dokładność,
    \item jakość inferencji online zależy nie tylko od modelu, ale także od jakości
          detekcji i preprocesingu obrazu.
\end{itemize}

\section{Ograniczenia i ryzyka}
\label{sec:ograniczenia_ryzyka}

Najważniejsze ograniczenia obecnej wersji:
\begin{enumerate}
    \item klasyfikacja oparta na pojedynczych klatkach bez modelowania sekwencji,
    \item brak obsługi znaków dynamicznych wymagających ruchu,
    \item podatność na zmiany oświetlenia i tła,
    \item brak kalibracji pod indywidualne różnice użytkowników.
\end{enumerate}

\section{Kierunki dalszego rozwoju}
\label{sec:dalszy_rozwoj}

Rozszerzenia, które mogą poprawić jakość i użyteczność systemu:
\begin{itemize}
    \item dodanie modelu sekwencyjnego (np. LSTM/Transformer) dla rozpoznawania
          gestów dynamicznych,
    \item wygładzanie predykcji w oknie czasowym,
    \item augmentacja danych i rozszerzenie zbioru treningowego o rzeczywiste
          próbki z kamery,
    \item budowa modułu oceny wydajności czasu inferencji na różnych urządzeniach.
\end{itemize}

\section{Wniosek końcowy}
\label{sec:wniosek_koncowy}

Projekt spełnia cele laboratorium z zakresu sztucznej inteligencji:
łączy techniki uczenia maszynowego i przetwarzania obrazu w działającym
prototypie klasyfikatora znaków dłoni, który może stanowić bazę do
dalszych prac badawczych i inżynierskich.
